\begin{conclusion}[分歧十次多项式的满对称伽罗瓦群判定]\label{con:fold-gauge-anomaly-p10-galois-s10}
设
\[
P_{10}(u)=27u^{10}+27u^9-153u^8-163u^7+793u^6+1061u^5-832u^4-816u^3+1320u^2-440u+40.
\]
则其分裂域在 \(\QQ\) 上的伽罗瓦群满足
\[
\mathrm{Gal}(P_{10}/\QQ)\cong S_{10}.
\]
并且
\[
\mathrm{Disc}(P_{10})=2^{38}\cdot 3^{24}\cdot 5^{2}\cdot 1931\cdot 9013^{2}\cdot 34319^{3}
\]
非平方，故该群不包含于 \(A_{10}\)。
\end{conclusion}
\begin{proof}[证明要点]
结论 \ref{con:fold-gauge-anomaly-disc-factorization-simple-branch} 已给出 \(P_{10}\) 的不可约性，故伽罗瓦群传递。
命题 \ref{prop:fold-gauge-anomaly-P10-galois-discriminant} 给出模 \(7\) 分解型 \((9,1)\)、模 \(11\) 分解型 \((7,3)\) 与模 \(13\) 不可约型 \((10)\)。
据 Dedekind--Frobenius 对应，群中含 \((9)(1)\) 与 \((7)(3)\) 型元素；由 \((7)(3)\) 的三次幂得到 \(7\)-循环。
传递子群含 \((9)(1)\) 可推出本原性；由 Jordan 定理（\(n=10,p=7\le n-3\)）得 \(A_{10}\subseteq G\)。
再由判别式非平方（命题 \ref{prop:fold-gauge-anomaly-P10-galois-discriminant}）排除 \(G\subseteq A_{10}\)，故 \(G=S_{10}\)。证毕。
\end{proof}

\begin{conclusion}[压力谱四次方程在 \texorpdfstring{$\QQ(u)$}{Q(u)} 上的泛型伽罗瓦群]\label{con:fold-gauge-anomaly-quartic-galois-s4}
令
\[
f(\mu,u):=\mu^4-\mu^3-(2u+1)\mu^2+(2u-u^3)\mu+2u.
\]
则在函数域 \(\QQ(u)\) 上有
\[
\mathrm{Gal}\bigl(f(\mu,u)/\QQ(u)\bigr)\cong S_4.
\]
\end{conclusion}
\begin{proof}[证明要点]
首先，\(u=2\) 特化给出四次多项式
\(
\mu^4-\mu^3-5\mu^2-4\mu+4
\)
在 \(\QQ\) 上不可约，从而泛型不可约。
其次，命题 \ref{prop:fold-gauge-anomaly-cubic-resolvent-discriminant} 给出
\(
\mathrm{Disc}_\mu(f)=-u(u-1)P_{10}(u)
\)
非平方，故伽罗瓦群不包含于 \(A_4\)。
再者，三次预解式
\[
\mathcal R(y,u)=y^3+(1+2u)y^2+(u^3-10u)y-(u^6-4u^4+20u^2+10u)
\]
在 \(u=2\) 处化为
\(
y^3+5y^2-12y-100
\)
且在 \(\QQ\) 上无有理根，故预解式泛型不可约。
四次群分类遂给出 \(S_4\)。
该结论亦可直接由结论 \ref{con:fold-gauge-anomaly-generic-galois-s4} 读取。证毕。
\end{proof}

\begin{conclusion}[谱四次覆盖与 \texorpdfstring{$S_4$}{S4} 伽罗瓦闭包的亏格公式]\label{con:fold-gauge-anomaly-genus-g3-g49}
设 \(\mathcal C_F\) 为谱四次 \(f(\mu,u)=0\) 的光滑模型，\(\pi:\mathcal C_F\to\PP^1_u\) 为投影。
则：
\begin{enumerate}
  \item \(\pi\) 为次数 \(4\) 覆盖，分支值恰为 \(u=0,1\) 与 \(P_{10}(u)=0\) 的十个根，共 \(12\) 个有限分支点，且全部简单分支；
  \item \(g(\mathcal C_F)=3\)；
  \item 其 \(S_4\)-伽罗瓦闭包 \(\mathcal X\to\PP^1_u\)（次数 \(24\)）满足 \(g(\mathcal X)=49\)。
\end{enumerate}
\end{conclusion}
\begin{proof}[证明要点]
第 (1) 条与第 (2) 条见命题 \ref{prop:fold-gauge-anomaly-spectral-quartic-geometry} 及结论 \ref{con:fold-gauge-anomaly-disc-factorization-simple-branch}。
对第 (3) 条，惯性均为换位（阶 \(2\)），每个分支点对 Riemann--Hurwitz 的贡献为 \(24(1-\tfrac12)=12\)。
故
\[
2g(\mathcal X)-2=24(-2)+12\cdot 12=96,
\]
从而 \(g(\mathcal X)=49\)。
亦与结论 \ref{con:fold-gauge-anomaly-s4-hurwitz-tower-genus} 一致。证毕。
\end{proof}

\begin{conclusion}[全纯微分的 Chevalley--Weil 精确分解]\label{con:fold-gauge-anomaly-chevalley-weil-s4}
在结论 \ref{con:fold-gauge-anomaly-genus-g3-g49} 的 \(\mathcal X\) 上，令 \(G=S_4\)。
则
\[
H^0(\mathcal X,\Omega^1)\cong
5\cdot \mathbf{sgn}\ \oplus\ 4\cdot \rho_{2}\ \oplus\ 3\cdot \rho_{3}\ \oplus\ 9\cdot \rho_{3'},
\]
其中 \(\mathbf{sgn}\) 为符号表示，\(\rho_2\) 为 \(2\) 维不可约表示，\(\rho_3\) 为标准 \(3\) 维表示，\(\rho_{3'}=\rho_3\otimes\mathbf{sgn}\)。
\end{conclusion}
\begin{proof}[证明要点]
分支点数 \(r=12\)，且全部惯性由换位生成。
Chevalley--Weil 公式给出
\[
m_\rho=-\dim\rho+\frac r2\bigl(\dim\rho-\dim\rho^\tau\bigr),
\]
其中 \(\tau\) 为换位。
由 \(S_4\) 角色表：
\(
\chi_{\mathbf{sgn}}(\tau)=-1,\chi_{\rho_2}(\tau)=0,\chi_{\rho_3}(\tau)=1,\chi_{\rho_{3'}}(\tau)=-1
\)，
故
\(
\dim(\mathbf{sgn})^\tau=0,\dim(\rho_2)^\tau=1,\dim(\rho_3)^\tau=2,\dim(\rho_{3'})^\tau=1
\)。
代入 \(r=12\) 得重数
\[
m_{\mathbf{sgn}}=5,\quad m_{\rho_2}=4,\quad m_{\rho_3}=3,\quad m_{\rho_{3'}}=9.
\]
维数核对：
\(
5\cdot1+4\cdot2+3\cdot3+9\cdot3=49=g(\mathcal X)
\)。
证毕。
\end{proof}

\begin{conclusion}[Jacobian 的 \texorpdfstring{$S_4$}{S4}-等变同源分解与 \texorpdfstring{$L$}{L}-函数因子化]\label{con:fold-gauge-anomaly-jacobian-equivariant-factorization}
设 \(J:=\mathrm{Jac}(\mathcal X)\)。
则存在定义于 \(\QQ\) 上的阿贝尔簇 \(A_{\mathbf{sgn}},A_2,A_3,A_{3'}\)，使
\[
J\sim_{\QQ}A_{\mathbf{sgn}}\times A_2\times A_3\times A_{3'}.
\]
其维数分别为
\[
\dim A_{\mathbf{sgn}}=5,\qquad \dim A_2=8,\qquad \dim A_3=9,\qquad \dim A_{3'}=27.
\]
并且在良约化素数处有 Euler 因子乘法分解，从而
\[
L(J,s)=L(A_{\mathbf{sgn}},s)\,L(A_2,s)\,L(A_3,s)\,L(A_{3'},s).
\]
\end{conclusion}
\begin{proof}[证明要点]
由结论 \ref{con:fold-gauge-anomaly-chevalley-weil-s4} 的表示分解，群代数中心幂等元在 \(J\) 上给出等变同源分解（Kani--Rosen 型机制）。
各分量维数由“表示维数 \(\times\) 重数”直接给出，即 \((5,8,9,27)\)。
\(\ell\)-进同调的半单分解与同源不变性给出 \(L\)-函数因子化。证毕。
\end{proof}

\begin{theorem}[\texorpdfstring{$S_4$}{S4} 正则闭包 Jacobian 的 N\'eron--Severi 秩下界]\label{thm:fold-gauge-anomaly-s4-jacobian-picard-rank-lower-bound-76}
令 \(J=\mathrm{Jac}(\mathcal X)\)，并在 \(\overline{\QQ}\) 上取其 N\'eron--Severi 群 \(\mathrm{NS}(J_{\overline{\QQ}})\)。
则
\[
\rho(J_{\overline{\QQ}}):=\mathrm{rank}\,\mathrm{NS}(J_{\overline{\QQ}})\ \ge\ 76.
\]
\leanverified{paper\_fold\_gauge\_anomaly\_s4\_jacobian\_picard\_rank\_lower\_bound\_76}
\end{theorem}
\begin{proof}
令
\(
V:=H^{0}(\mathcal X,\Omega^{1})
\)。
由结论 \ref{con:fold-gauge-anomaly-chevalley-weil-s4}，
\[
V\cong
5\cdot \mathbf{sgn}\ \oplus\ 4\cdot \rho_{2}\ \oplus\ 3\cdot \rho_{3}\ \oplus\ 9\cdot \rho_{3'}.
\]
由于 \(S_4\) 的不可约表示均可在 \(\QQ\) 上实现，且对每个不可约 \(\rho\) 有 \(\operatorname{End}_{S_4}(\rho)\cong\QQ\)，
故其交换子代数满足
\[
\operatorname{End}_{S_4}(V)\ \cong\ M_{5}(\QQ)\times M_{4}(\QQ)\times M_{3}(\QQ)\times M_{9}(\QQ).
\]
\par
另一方面，\(S_4\) 作用于 \(\mathcal X\) 诱导 \(\QQ[S_4]\) 在 \(J\) 上的对应作用，并在
\(
H^{0}(\mathcal X,\Omega^{1})
\)
上给出与上式兼容的表示；因此 \(\operatorname{End}_{S_4}(V)\) 可视为 \(\operatorname{End}^{0}(J)\) 中的一族 \(\QQ\)-子代数（在同源范畴中由等变分解的重数空间给出矩阵单元）。
\par
取 \(J\) 的主极化所诱导的 Rosati 反自同态 \(\dagger\)。
在上述分解下，\(\dagger\) 在每个矩阵因子 \(M_m(\QQ)\) 上诱导转置反演；因此 Rosati 不变子空间维数至少为
\[
\sum_{m\in\{5,4,3,9\}}\dim_{\QQ}\mathrm{Sym}_{m}(\QQ)
\;=\;
\sum_{m\in\{5,4,3,9\}}\frac{m(m+1)}{2}
\;=\;76.
\]
最后，对任意主极化阿贝尔簇 \(A\) 有标准同构
\[
\mathrm{NS}(A)_{\QQ}\ \cong\ \operatorname{End}^{0}(A)^{\dagger=1},
\]
从而 \(\rho(J_{\overline{\QQ}})=\dim_{\QQ}\mathrm{NS}(J_{\overline{\QQ}})_{\QQ}\ge 76\)。证毕。
\end{proof}

\begin{theorem}[换位分歧下的驯 Artin 导数与函数域导子度]\label{thm:fold-gauge-anomaly-s4-regular-closure-tame-artin-conductor-degrees}
沿用结论 \ref{con:fold-gauge-anomaly-genus-g3-g49} 的正则闭包覆盖
\(\mathcal X\to\PP^1_u\)。
其有限分歧点数为 \(12\)，且每个惰性群由一换位 \(\tau\in S_4\) 生成。
对 \(S_4\) 的四个非平凡不可约表示 \(\rho\in\{\mathbf{sgn},\rho_2,\rho_3,\rho_{3'}\}\)，定义驯 Artin 导子指数
\[
a(\rho;\tau):=\dim(\rho)-\dim\bigl(\rho^{\langle\tau\rangle}\bigr).
\]
则
\[
a(\mathbf{sgn};\tau)=1,\qquad
a(\rho_{2};\tau)=1,\qquad
a(\rho_{3};\tau)=1,\qquad
a(\rho_{3'};\tau)=2.
\]
从而相应的全局导子度（导子除子之次数）
\(
A(\rho):=\sum_{b}a(\rho;I_b)
\)
满足
\[
A(\mathbf{sgn})=12,\qquad
A(\rho_{2})=12,\qquad
A(\rho_{3})=12,\qquad
A(\rho_{3'})=24,
\]
其中求和遍及全部 \(12\) 个分歧点。
\leanverified{paper\_fold\_gauge\_anomaly\_s4\_artin\_conductor\_degrees}
\end{theorem}
\begin{proof}
由于覆盖在每个分歧点处为驯分歧且惰性群为阶 \(2\) 的循环群 \(\langle\tau\rangle\)，有
\(
a(\rho;\tau)=\dim\rho-\dim\rho^{\langle\tau\rangle}
\)。
而对阶 \(2\) 作用，\(\rho^{\langle\tau\rangle}\) 的维数由角色平均给出：
\[
\dim\rho^{\langle\tau\rangle}=\frac{\dim\rho+\operatorname{tr}(\rho(\tau))}{2}.
\]
对 \(S_4\) 的换位 \(\tau\)，角色值为
\[
\operatorname{tr}(\mathbf{sgn}(\tau))=-1,\qquad
\operatorname{tr}(\rho_2(\tau))=0,\qquad
\operatorname{tr}(\rho_3(\tau))=1,\qquad
\operatorname{tr}(\rho_{3'}(\tau))=-1.
\]
代入得到
\[
\dim(\mathbf{sgn})^{\langle\tau\rangle}=0,\quad
\dim(\rho_2)^{\langle\tau\rangle}=1,\quad
\dim(\rho_3)^{\langle\tau\rangle}=2,\quad
\dim(\rho_{3'})^{\langle\tau\rangle}=1,
\]
从而 \(a(\rho;\tau)\) 如陈述所示。
由于共有 \(12\) 个分歧点且局部数据一致，故 \(A(\rho)=12\cdot a(\rho;\tau)\)，得到全局导子度。证毕。
\end{proof}

\leanverified{paper\_fold\_gauge\_anomaly\_p10\_unique\_quadratic\_subfield}
\begin{corollary}[分裂域的唯一二次子域]\label{cor:fold-gauge-anomaly-p10-unique-quadratic-subfield}
设 \(L\) 为 \(P_{10}\) 的分裂域。
则 \(L\) 存在且仅存在一个二次子域，且
\[
L^{A_{10}}=\QQ\!\bigl(\sqrt{\mathrm{Disc}(P_{10})}\bigr)=\QQ\!\bigl(\sqrt{66269989}\bigr).
\]
\end{corollary}
\begin{proof}
由结论 \ref{con:fold-gauge-anomaly-p10-galois-s10}，\(\mathrm{Gal}(L/\QQ)\cong S_{10}\)。
\(S_{10}\) 的唯一指数 \(2\) 正规子群为 \(A_{10}\)，故二次子域唯一。
判别式平方自由部分为 \(1931\cdot 34319=66269989\)，即得所述表达。证毕。
\end{proof}

\begin{corollary}[分裂域次数与分支根不可根式性]\label{cor:fold-gauge-anomaly-p10-degree-and-unsolvability}
对 \(P_{10}\) 的分裂域 \(L\) 有
\[
[L:\QQ]=|S_{10}|=10!=3{,}628{,}800.
\]
从而 \(P_{10}(u)\) 的任一根均不可由根式表达。
\end{corollary}
\begin{proof}
由结论 \ref{con:fold-gauge-anomaly-p10-galois-s10}，分裂域次数等于伽罗瓦群阶，故为 \(10!\)。
又 \(S_{10}\) 非可解，依 Galois--Abel 判别，根式解不存在。证毕。
\leanverified{paper\_fold\_gauge\_anomaly\_p10\_degree\_and\_unsolvability}
\end{proof}

\begin{corollary}[两个正规商曲线的 Jacobian 等型实现]\label{cor:fold-gauge-anomaly-quotient-jacobians-isotypic}
设
\[
\mathcal C_{\mathrm{disc}}:=\mathcal X/A_4=\mathcal H,
\qquad
\mathcal Y_{V_4}:=\mathcal X/V_4.
\]
则有同源分解
\[
\mathrm{Jac}(\mathcal C_{\mathrm{disc}})\sim_{\QQ}A_{\mathbf{sgn}},
\qquad
\mathrm{Jac}(\mathcal Y_{V_4})\sim_{\QQ}A_{\mathbf{sgn}}\times A_2.
\]
\end{corollary}
\begin{proof}
对任意正规子群 \(N\triangleleft S_4\)，有
\[
H^0(\mathcal X/N,\Omega^1)\cong H^0(\mathcal X,\Omega^1)^N.
\]
由结论 \ref{con:fold-gauge-anomaly-chevalley-weil-s4}，
\[
H^0(\mathcal X,\Omega^1)\cong
5\cdot \mathbf{sgn}\oplus 4\cdot \rho_2\oplus 3\cdot \rho_3\oplus 9\cdot \rho_{3'}.
\]
对 \(A_4\) 取不变向量时，仅 \(\mathbf{sgn}\) 在 \(A_4\) 上平凡，故
\(
H^0(\mathcal X,\Omega^1)^{A_4}\cong 5\cdot\mathbf{sgn}
\)，
从而 \(\mathrm{Jac}(\mathcal X/A_4)\sim_{\QQ}A_{\mathbf{sgn}}\)。
对 \(V_4\) 取不变向量时，\(\mathbf{sgn}\) 与 \(\rho_2\) 平凡而 \(\rho_3,\rho_{3'}\) 非平凡，故
\(
H^0(\mathcal X,\Omega^1)^{V_4}\cong 5\cdot\mathbf{sgn}\oplus 4\cdot\rho_2
\)，
进而
\(
\mathrm{Jac}(\mathcal X/V_4)\sim_{\QQ}A_{\mathbf{sgn}}\times A_2
\)。
证毕。
\leanverified{paper\_fold\_gauge\_anomaly\_quotient\_jacobians\_isotypic}
\end{proof}

\begin{theorem}[\texorpdfstring{$A_2$}{A2} 因子的 Prym 实现]\label{thm:fold-gauge-anomaly-a2-prym-realization}
由 \(V_4\triangleleft A_4\triangleleft S_4\) 得到自然覆盖
\[
\pi:\mathcal Y_{V_4}=\mathcal X/V_4\longrightarrow \mathcal C_{\mathrm{disc}}=\mathcal X/A_4,
\qquad
\mathrm{Gal}(\pi)\cong A_4/V_4\cong C_3.
\]
其 Prym 簇
\[
\operatorname{Prym}(\pi):=\ker\!\bigl(\operatorname{Nm}_{\pi}:\mathrm{Jac}(\mathcal Y_{V_4})\to \mathrm{Jac}(\mathcal C_{\mathrm{disc}})\bigr)^0
\]
满足
\[
\operatorname{Prym}(\pi)\sim_{\QQ}A_2,
\qquad
\dim \operatorname{Prym}(\pi)=8.
\]
\end{theorem}
\begin{proof}
在结论 \ref{con:fold-gauge-anomaly-disc-factorization-simple-branch} 的 \(S_4\)-闭包中，有限分支点的惯性均由换位生成；
换位属于 \(S_4\setminus A_4\)，故其与 \(A_4\)（进而与 \(V_4\)）的交均为平凡。
因此 \(\pi\) 为循环三次非分歧覆盖。
对非分歧循环三次覆盖，Jacobian 有标准分解
\[
\mathrm{Jac}(\mathcal Y_{V_4})\sim_{\QQ}\mathrm{Jac}(\mathcal C_{\mathrm{disc}})\times \operatorname{Prym}(\pi).
\]
再由推论 \ref{cor:fold-gauge-anomaly-quotient-jacobians-isotypic}，
\(
\mathrm{Jac}(\mathcal Y_{V_4})\sim_{\QQ}A_{\mathbf{sgn}}\times A_2
\)
且
\(
\mathrm{Jac}(\mathcal C_{\mathrm{disc}})\sim_{\QQ}A_{\mathbf{sgn}}
\)，
比较即得 \(\operatorname{Prym}(\pi)\sim_{\QQ}A_2\)。
维数由
\(
g(\mathcal Y_{V_4})-g(\mathcal C_{\mathrm{disc}})=13-5=8
\)
给出。证毕。
\end{proof}

\leanverified{paper\_fold\_gauge\_anomaly\_a2\_prym\_realization}

\leanverified{paper\_fold\_gauge\_anomaly\_a2\_eisenstein\_multiplication}
\begin{theorem}[\texorpdfstring{$A_2$}{A2} 的 Eisenstein 整数乘法]\label{thm:fold-gauge-anomaly-a2-eisenstein-multiplication}
在定理 \ref{thm:fold-gauge-anomaly-a2-prym-realization} 的记号下，存在环嵌入
\[
\ZZ[\zeta_3]\hookrightarrow \operatorname{End}\!\bigl(\operatorname{Prym}(\pi)\bigr),
\qquad
\QQ(\zeta_3)\hookrightarrow \operatorname{End}^0\!\bigl(\operatorname{Prym}(\pi)\bigr).
\]
等价地，
\[
\QQ(\sqrt{-3})\subseteq \operatorname{End}^0(A_2).
\]
\end{theorem}
\begin{proof}
令 \(\sigma\) 为 \(\pi\) 的阶 \(3\) 甲板变换生成元。
\(\sigma\) 作用于 \(\mathrm{Jac}(\mathcal Y_{V_4})\)，并与范数映射 \(\operatorname{Nm}_{\pi}=1+\sigma+\sigma^2\) 交换，故保持
\(\operatorname{Prym}(\pi)=\ker(\operatorname{Nm}_{\pi})^0\)。
在 Prym 上，关系 \(\sigma^3=1\) 与 \(1+\sigma+\sigma^2=0\) 同时成立，故
\(\sigma\) 满足最小多项式 \(x^2+x+1\)，从而诱导
\(
\ZZ[\zeta_3]\cong \ZZ[\sigma]\subseteq \operatorname{End}(\operatorname{Prym}(\pi))
\)。
张量 \(\QQ\) 即得 \(\QQ(\zeta_3)\subseteq \operatorname{End}^0(\operatorname{Prym}(\pi))\)。
再由定理 \ref{thm:fold-gauge-anomaly-a2-prym-realization} 的同源识别
\(\operatorname{Prym}(\pi)\sim_{\QQ}A_2\) 得到结论。证毕。
\end{proof}

\begin{conjecture}[等型因子的绝对单纯与内自同态环极小化]\label{conj:fold-gauge-anomaly-isotypic-absolute-simplicity}
在结论 \ref{con:fold-gauge-anomaly-jacobian-equivariant-factorization} 的分解
\[
\mathrm{Jac}(\mathcal X)\sim_{\QQ}A_{\mathbf{sgn}}\times A_2\times A_3\times A_{3'}
\]
中，预期成立：
\begin{enumerate}
  \item \(A_2\) 在 \(\overline{\QQ}\) 上绝对单纯，且
  \[
  \operatorname{End}^0_{\overline{\QQ}}(A_2)=\QQ(\zeta_3).
  \]
  \item \(A_3,A_{3'}\) 在 \(\overline{\QQ}\) 上绝对单纯，且
  \[
  \operatorname{End}^0_{\overline{\QQ}}(A_3)=\operatorname{End}^0_{\overline{\QQ}}(A_{3'})=\QQ.
  \]
\end{enumerate}
该预期与通用简单分歧 \(S_4\)-覆盖的 Mumford--Tate 极大性一致；其可通过多组良素数 Frobenius 特征多项式与 Honda--Tate 判别进行有效核验。
\end{conjecture}

\endinput
